%%%%%%%%%%%%%%%%%%%%%%%%%%%%%%%%%%%%%%%%%%%%%%%%%%%
%% Exercices & Solutions — Informatique Quantique
%% Stage LIA / CERI — Université d'Avignon
%% Stagiaire : Patrice SEBASTIANO
%% Encadrant : Serigne GUEYE
%%%%%%%%%%%%%%%%%%%%%%%%%%%%%%%%%%%%%%%%%%%%%%%%%%%

% \documentclass[light]{ceri/sty/rapport}   % compilation rapide
% \documentclass[handout]{ceri/sty/rapport}  % impression
% \documentclass[full]{ceri/sty/rapport}     % avec listes des figures/tableaux
\documentclass{ceri/sty/rapport}

%%%%%%%%%%%%%%%%%%%%%%%%%%%%%%%%%%%%%%%%%%%%%%%%%%%
%% Informations générales
%%%%%%%%%%%%%%%%%%%%%%%%%%%%%%%%%%%%%%%%%%%%%%%%%%%
\major{Master d'Informatique}
\specialization{Informatique Quantique}
\course{Stage de Recherche — LIA Avignon}

\title{Exercices \& Solutions}
\subtitle{Informatique Quantique}

\author{
Patrice SEBASTIANO
}

\advisor[Encadrant]{
Serigne GUEYE
}

%%%%%%%%%%%%%%%%%%%%%%%%%%%%%%%%%%%%%%%%%%%%%%%%%%%
%% Packages supplémentaires
%%%%%%%%%%%%%%%%%%%%%%%%%%%%%%%%%%%%%%%%%%%%%%%%%%%
\usepackage{ifthen}
\usepackage{amsmath,amssymb}
\usepackage{mathtools}
\usepackage{braket}
\usepackage{quantikz}
\usepackage{tcolorbox}
\tcbuselibrary{breakable, skins}

% Couleurs personnalisées
\definecolor{exo-bg}{RGB}{230, 240, 255}
\definecolor{exo-border}{RGB}{50, 100, 200}
\definecolor{sol-bg}{RGB}{230, 255, 235}
\definecolor{sol-border}{RGB}{40, 160, 60}

% Environnement exercice — argument optionnel = référence page du livre
\newcounter{exercice}[section]
\renewcommand{\theexercice}{\thesection.\arabic{exercice}}

\newtcolorbox[use counter=exercice]{exercice}[1][]{%
colback=exo-bg,
colframe=exo-border,
fonttitle=\bfseries,
title={Exercice~\theexercice \ifthenelse{\equal{#1}{}}{}{~—~#1}},
breakable,
before upper={\parindent=1em}
}

% Environnement solution — argument optionnel transmis à tcolorbox
\newtcolorbox{solution}[1][]{%
colback=sol-bg,
colframe=sol-border,
fonttitle=\bfseries,
title={Solution},
breakable,
before upper={\parindent=1em},
#1
}

% Environnement démonstration
\definecolor{demo-bg}{RGB}{248, 240, 255}
\definecolor{demo-border}{RGB}{120, 50, 180}

\newtcolorbox{demonstration}[1][]{%
colback=demo-bg,
colframe=demo-border,
fonttitle=\bfseries,
title={Démonstration\ifthenelse{\equal{#1}{}}{}{~—~#1}},
enhanced jigsaw,
breakable,
before upper={\parindent=1em}
}

%%%%%%%%%%%%%%%%%%%%%%%%%%%%%%%%%%%%%%%%%%%%%%%%%%%
%% Début du document
%%%%%%%%%%%%%%%%%%%%%%%%%%%%%%%%%%%%%%%%%%%%%%%%%%%
\begin{document}

% Suppress the class's internal "Titre" and "Sommaire" TOC entries
\let\origaddcontentsline\addcontentsline
\renewcommand{\addcontentsline}[3]{}
\maketitle
\let\addcontentsline\origaddcontentsline
\sloppy

\phantomsection
\addcontentsline{toc}{section}{Démonstrations}
{\color{fgRed}\normalfont\Large\bfseries Démonstrations}
\par\vspace{2mm}
\begin{demonstration}[Addition scalaire dans la base de Fourier par paroi de portes de phase]

\textbf{Contexte.}
Ce résultat (p.~229--230 du livre de référence) justifie pourquoi une
\emph{paroi de portes de phase} appliquée à un registre encodé en base QFT
réalise une addition scalaire dans cette base.

\medskip
\textbf{Théorème.}
Soit $|\psi_j\rangle$ l'encodage QFT d'un entier $j$, et soit
$\mathcal{P} = \bigotimes_{h=0}^{m-1} P\!\bigl(\frac{\pi l}{2^h}\bigr)$
une paroi de portes de phase paramétrée par un entier $l$. Alors :
\[
  \mathcal{P}\,|\psi_j\rangle = |\psi_{j+l}\rangle.
\]
Appliquer $\mathcal{P}$ \emph{additionne} $l$ à la valeur encodée dans la
base de Fourier.

\bigskip

Soit $m$ le nombre de qubits, indexés de $h = 0$ (MSB) à $h = m-1$ (LSB).
Un état de base est $|\underline{k}\rangle = |k_0 k_1 \dots k_{m-1}\rangle$
avec $k = \sum_{h=0}^{m-1} k_h\, 2^{m-1-h}$.

\medskip
\noindent\textbf{Étape 1 — Action de $\mathcal{P}$ sur un état de base $|\underline{k}\rangle$.}

\smallskip
Par distributivité du produit tensoriel :
\begin{align*}
\mathcal{P}|\underline{k}\rangle
  &= \bigotimes_{h=0}^{m-1} P\!\left(\frac{\pi l}{2^h}\right) |k_h\rangle \\
  &= \prod_{\substack{h=0 \\ k_h \neq 0}}^{m-1}
       \exp\!\left(\frac{i\pi l}{2^h}\right)
       \cdot \bigotimes_{h=0}^{m-1} |k_h\rangle \\
  &= \exp\!\left(\sum_{h=0}^{m-1} \frac{i\pi l\, k_h}{2^h}\right) |\underline{k}\rangle.
\end{align*}

En mettant $2^{m-1}$ au dénominateur commun :
\begin{align*}
\sum_{h=0}^{m-1} \frac{i\pi l\, k_h}{2^h}
  &= i\pi l \cdot
     \frac{\displaystyle\sum_{h=0}^{m-1} k_h\, 2^{m-1-h}}{2^{m-1}}
   = \frac{i\pi l\cdot k}{2^{m-1}}
   = \frac{2\pi i\, l\cdot k}{2^m}.
\end{align*}

On obtient donc :
\[
  \mathcal{P}|\underline{k}\rangle
  = \exp\!\left(\frac{2\pi i\, l\cdot k}{2^m}\right)|\underline{k}\rangle.
\]

\medskip
\noindent\textbf{Étape 2 — Addition sur un état encodé en base QFT.}

\smallskip
L'encodage QFT de $j$ est :
\[
  |\psi_j\rangle
  = \frac{1}{\sqrt{2^m}} \sum_{k=0}^{2^m-1}
    \exp\!\left(\frac{2\pi i\, j\cdot k}{2^m}\right) |\underline{k}\rangle.
\]

Par linéarité et en utilisant le résultat de l'étape 1 :
\begin{align*}
\mathcal{P}|\psi_j\rangle
  &= \frac{1}{\sqrt{2^m}} \sum_{k=0}^{2^m-1}
       \exp\!\left(\frac{2\pi i\, j\cdot k}{2^m}\right)
       \exp\!\left(\frac{2\pi i\, l\cdot k}{2^m}\right) |\underline{k}\rangle \\
  &= \frac{1}{\sqrt{2^m}} \sum_{k=0}^{2^m-1}
       \exp\!\left(\frac{2\pi i\,(j+l)\cdot k}{2^m}\right) |\underline{k}\rangle
   = |\psi_{j+l}\rangle.
\end{align*}

\[
  \square
\]

\end{demonstration}

\bigskip

\setcounter{section}{5}
\section{Exercices du Chapitre 6}
\label{sec:chap6}

\begin{exercice}
Prove that $O_f$ is not only reversible, but also unitary, and hence deserves the name “quantum gate”.

The context is Grover search: the oracle $O_f$ acts as
\[ O_f\ket{x}\ket{y} = \ket{x}\ket{y\oplus f(x)}. \]
\end{exercice}

\begin{solution}
Pour la reversibilite, on remarque que l'application de $O_f$ deux fois donne de nouveau l'etat d'origine :
\[
O_f\bigl(O_f\ket{x}\ket{y}\bigr)
= O_f\ket{x}\ket{y\oplus f(x)}
= \ket{x}\ket{(y\oplus f(x))\oplus f(x)}
= \ket{x}\ket{y}.
\]
Ainsi $O_f^2 = I$, ce qui est equivalent a dire que $O_f$ est reversible.

Autre raisonnement pour la reversibilite : en sortie de l'oracle on connait $x$, donc on connait $f(x)$. On connait egalement $y\oplus f(x)$ en second registre. Comme $f(x)$ est connu, on peut retrouver $y$ en effectuant a nouveau l'operation $\oplus f(x)$ sur ce second registre. Donc l'operation est reversible.

Pour l'unitarite, on considere $O_f$ comme une matrice a controle multiple : les qubits de $x$ servent a controler le bit $y$. Comme $x$ est encode sur $n$ qubits, il existe $2^n$ configurations possibles pour $x$.

Pour chaque configuration $x$ on a soit $f(x)=1$, soit $f(x)=0$. Si $f(x)=1$, le registre $y$ est inverse par un operateur $X_2$. Si $f(x)=0$, le registre $y$ reste unchanged par l'operateur $I_2$.

Donc $O_f$ est une matrice par blocs, avec $2^n$ blocs diagonaux de dimension $2$ chacun. Ces blocs sont soit $X$, soit $I$.

Comme $O_f$ est bloc-diagonale, son adjoint $O_f^\dagger$ se calcule bloc par bloc. Or $I$ et $X$ sont unitaires, donc $I^\dagger = I$ et $X^\dagger = X$. Par consequent $O_f^\dagger = O_f$ et on obtient
\[
O_f^\dagger O_f = I.
\]
Cela montre que $O_f$ est unitaire.

Ainsi $O_f$ est a la fois reversible et unitaire, ce qui justifie pleinement l'appellation de "quantum gate".
\end{solution}

\begin{exercice}
Construct a quantum circuit for $O_f$ where $f$ is a 4-bit Boolean function that takes the value $1$ on $0111$, $1110$, and $0101$, and takes the value $0$ on every other input.
\end{exercice}

\begin{solution}
A circuit implementing $O_f$ can be built as the XOR of three multi-controlled NOT operations, one for each input where $f(x)=1$.
For $0111$, $1110$, and $0101$, the circuit is:
\[
\begin{quantikz}
\lstick{$q_0$} & \gate{X} & \ctrl{4} & \gate{X} & \qw      & \ctrl{4} & \qw      & \gate{X} & \ctrl{4} & \gate{X} & \qw \\
\lstick{$q_1$} & \qw      & \ctrl{3} & \qw      & \qw      & \ctrl{3} & \qw      & \qw      & \ctrl{3} & \qw      & \qw \\
\lstick{$q_2$} & \qw      & \ctrl{2} & \qw      & \qw      & \ctrl{2} & \qw      & \gate{X} & \ctrl{2} & \gate{X} & \qw \\
\lstick{$q_3$} & \qw      & \ctrl{1} & \qw      & \gate{X} & \ctrl{1} & \gate{X} & \qw      & \ctrl{1} & \qw      & \qw \\
\lstick{$y$}   & \qw      & \targ{}  & \qw      & \qw      & \targ{}  & \qw      & \qw      & \targ{}  & \qw      & \qw
\end{quantikz}
\]

The first multi-controlled NOT selects $0111$ by inverting $q_0$ before and after the control.
The second selects $1110$ by inverting $q_3$ before and after the control.
The third selects $0101$ by inverting $q_0$ and $q_2$ before and after the control.
Each target is the ancilla qubit $y$, so the oracle flips $y$ exactly when $f(x)=1$.
\end{solution}

\begin{exercice}
Using two's complement with 5 qubits, represent $10$ and $-7$ and perform their addition.
\end{exercice}

\begin{solution}

\textbf{Standard binary encoding — two useful patterns.}
In base-2, any $n$-bit string has positional values:
\[
  \underbrace{11\cdots1}_{n} = 2^n - 1,
  \qquad
  1\underbrace{0\cdots0}_{n-1} = 2^n - 2^{n-1}.
\]
These identities simplify decomposing a decimal integer into its binary representation.

\medskip\noindent\textbf{Two's complement convention.}
Integers are represented modulo $2^n$, so by definition:
\[
  -x \;\equiv\; 2^n - x \pmod{2^n}.
\]
Under this convention, the most significant bit (MSB) acts as a sign bit: $0$ for non-negative values, $1$ for negative values.
The register still encodes exactly $2^n$ distinct values, but over the interval $[-2^{n-1},\, 2^{n-1}-1]$ instead of $[0,\, 2^n-1]$.

\medskip\noindent\textbf{Encoding $+10$.}
\[
  10 = 8 + 2 = \underbrace{1000}_{2^3} + \underbrace{0010}_{2^1}
     = \underbrace{0}_{s}1010.
\]
The sign bit is $0$ (positive).

\medskip\noindent\textbf{Encoding $-7$.}
\[
  -7 \;\equiv\; 2^5 - 7 = 32 - 7 = 25.
\]
\[
  25 = 16 + 8 + 1
     = \underbrace{10000}_{2^4} + \underbrace{01000}_{2^3} + \underbrace{00001}_{2^0}
     = 11001.
\]
The sign bit is $1$, confirming a negative value.

\medskip\noindent\textbf{Addition.}
\[
\begin{array}{r}
  01010 \quad (+10) \\
+\; 11001 \quad (-7) \\\hline
\mathbf{1}\,00011
\end{array}
\]
The leading carry $\mathbf{1}$ overflows beyond the 5-qubit register and is discarded (5-bit modular arithmetic).

\medskip\noindent\textbf{Result.}
The 5-bit register holds $00011 = 3$.
\[
  10 + (-7) = 3. \quad\checkmark
\]

\end{solution}

\begin{exercice}
Derive a circuit that prepares the \emph{phase representation} of $0$, adds $6$
to it, and then subtracts $4$. Use $4$ qubits.
\end{exercice}

\begin{solution}
\medskip\noindent\textbf{Arithmetic in the phase basis.}
Adding a constant $c$ multiplies the phase of qubit $k$ by $e^{2\pi ic/2^k}$.
The single-qubit phase gate $P(\theta)=\mathrm{diag}(1,e^{i\theta})$ achieves
this: applying $P\!\bigl(\tfrac{2\pi c}{2^k}\bigr)$ to qubit $k$ shifts its
phase by the required amount.
For $c=+6$ the angles are $6\pi,\,3\pi,\,\tfrac{3\pi}{2},\,\tfrac{3\pi}{4}$
(qubits $1$–$4$); for $c=-4$ they are
$-4\pi,\,-2\pi,\,-\pi,\,-\tfrac{\pi}{2}$.

\medskip\noindent\textbf{Circuit.}
\[
\begin{quantikz}
\lstick{$\ket{0}$} & \gate{H} & \gate{\ensuremath{P(6\pi)}}                        & \gate{\ensuremath{P(-4\pi)}}                        & \qw \\
\lstick{$\ket{0}$} & \gate{H} & \gate{\ensuremath{P\!\bigl(\tfrac{6\pi}{2}\bigr)}} & \gate{\ensuremath{P\!\bigl(\tfrac{-4\pi}{2}\bigr)}} & \qw \\
\lstick{$\ket{0}$} & \gate{H} & \gate{\ensuremath{P\!\bigl(\tfrac{6\pi}{4}\bigr)}} & \gate{\ensuremath{P\!\bigl(\tfrac{-4\pi}{4}\bigr)}} & \qw \\
\lstick{$\ket{0}$} & \gate{H} & \gate{\ensuremath{P\!\bigl(\tfrac{6\pi}{8}\bigr)}} & \gate{\ensuremath{P\!\bigl(\tfrac{-4\pi}{8}\bigr)}} & \qw
\end{quantikz}
\]

\end{solution}

\begin{exercice}
Derive a circuit in the phase (QFT) basis for computing
$f(x_0, x_1, x_2) = x_1 x_2 - 3x_0 + 2$,
where $x_0, x_1, x_2 \in \{0,1\}$.
Determine the minimum output register size $m$ using the two's-complement
convention, then construct the full circuit.
\end{exercice}

\begin{solution}
\textbf{Step 1 — minimum register size.}
Define
\[
  K = \max\!\left(\sum_{g_i > 0} g_i,\; \sum_{g_j < 0} |g_j|\right)
    = \max(1 + 2,\; 3) = 3.
\]
A two's-complement register of $m$ bits covers $[-2^{m-1},\, 2^{m-1}-1]$.
The condition $K \le 2^{m-1}-1$ gives $3 \le 2^{m-1}-1$, hence $\mathbf{m = 3}$
(range $[-4,\, 3]$). $\checkmark$

\medskip\noindent\textbf{Step 2 — circuit design.}
The three output qubits $q_1, q_2, q_3$ (top to bottom, $k = 1, 2, 3$) are
prepared in the phase representation of $0$ via $H^{\otimes 3}$.
Each arithmetic term is added by phase rotations $P\!\bigl(\tfrac{2\pi c}{2^k}\bigr)$:

\begin{itemize}
  \item \textbf{Constant $+2$}: unconditional gates on every output qubit.
  \item \textbf{Term $+x_1 x_2$}: doubly-controlled gates (controls $x_1$ \emph{and} $x_2$) on every output qubit.
  \item \textbf{Term $-3x_0$}: singly-controlled gates (control $x_0$) on every output qubit.
\end{itemize}

\medskip\noindent\textbf{Circuit.}
\noindent\makebox[\linewidth][c]{%
\scalebox{0.78}{%
\begin{quantikz}[column sep=3pt]
\lstick{$x_0$}     & \qw\slice{\ensuremath{+2}}           & \qw\slice{\ensuremath{+x_1x_2}}     & \qw                        & \qw                          & \qw\slice{\ensuremath{-3x_0}}        & \ctrl{3}                     & \ctrl{4}                       & \ctrl{5}                       & \qw \\
\lstick{$x_1$}     & \qw                                  & \qw                                  & \ctrl{2}                   & \ctrl{3}                     & \ctrl{4}                             & \qw                          & \qw                            & \qw                            & \qw \\
\lstick{$x_2$}     & \qw                                  & \qw                                  & \ctrl{1}                   & \ctrl{2}                     & \ctrl{3}                             & \qw                          & \qw                            & \qw                            & \qw \\
\lstick{$\ket{0}$} & \gate{H}                             & \gate{\ensuremath{P(2\pi)}}          & \gate{\ensuremath{P(\pi)}} & \qw                          & \qw                                  & \gate{\ensuremath{P(-3\pi)}} & \qw                            & \qw                            & \qw \\
\lstick{$\ket{0}$} & \gate{H}                             & \gate{\ensuremath{P(\pi)}}           & \qw                        & \gate{\ensuremath{P(\pi/2)}} & \qw                                  & \qw                          & \gate{\ensuremath{P(-3\pi/2)}} & \qw                            & \qw \\
\lstick{$\ket{0}$} & \gate{H}                             & \gate{\ensuremath{P(\pi/2)}}         & \qw                        & \qw                          & \gate{\ensuremath{P(\pi/4)}}         & \qw                          & \qw                            & \gate{\ensuremath{P(-3\pi/4)}} & \qw
\end{quantikz}%
}%
}

\noindent\textit{Note:} $P(2\pi) = I$, so the gate on $q_1$ in the $+2$ stage is
the identity and may be omitted.
Applying the inverse QFT (IQFT) to the output register after the circuit recovers
the two's-complement binary encoding of $f(x_0, x_1, x_2)$.
\end{solution}

\end{document}
